\chapter{The conformal frame and $\Omega$}\label{ch:conformal}

\begin{questionbox}
The cap and the neck look very different: one is a shrinking sphere, the other an endless tube. Is there a frame in which they look alike? Yes. Both are conformally the Einstein static universe $\mathbb R\times S^3$, multiplied by a single function $\Omega$, and in that frame Einstein's equations collapse to the linear equation $\hat\Box\Omega=\Omega$. This chapter also contains the linear rigidity theorem, which the author believes to be new.
\end{questionbox}

\section{Why a conformal frame}

Yang--Mills theory in four dimensions does not care about the overall scale of the metric. If $g=\Omega^2\hat g$, a Yang--Mills field on $g$ is also one on $\hat g$, and its stress is traceless. Choose $\hat g=-d\eta^2+du^2+\sin^2u\,d\Omega^2$, the static universe whose space is the unit three-sphere: the quantum sphere of Chapter~6, now as space. The trace of Einstein's equations says $R[g]=0$. The conformal transformation formula for the scalar curvature (quoted \external) then turns this into
\[
\hat\Box\,\Omega=\Omega ,
\]
a linear wave equation for the conformal factor. That is a remarkable simplification for a non-linear theory.

\section{Cap and neck as conformal factors}

\begin{theorem}\proved
The cap is $\Omega=\ell\cos\eta$ (depending only on time) and the neck is $\Omega=\ell/\sin u$ (depending only on position). Both satisfy $\hat\Box\Omega=\Omega$. On both, the $\eta\eta$ component of Einstein's equations reduces to $\ell^2=4\pi G/g^2$.
\end{theorem}
\noindent In the repository: \lean{cap\_box}, \lean{neck\_box}, \lean{eta\_cap}, \lean{eta\_neck}.

The neck's coordinate $x=\ell\log\tan(u/2)$ turns $(\ell^2/\sin^2u)(du^2+\sin^2u\,d\Omega^2)$ into $dx^2+\ell^2d\Omega^2$ (\lean{neck\_coordinate}): the two infinitely distant ends of the neck are the two poles of $S^3$. At the junction $u=\pi/2$ both factors equal $\ell$ with zero slope, and the second derivatives are $0$ and $\ell$ (\lean{junction\_seam}). The $C^2$ seam of Chapter~16 is a seam in $\Omega''$.

The areal radius is $r=\Omega\sin\chi$, and it satisfies the flat two-dimensional wave equation (\lean{oq\_radius\_wave}). So $r=F(\eta+\chi)-F(\eta-\chi)$ for some function $F$, and the antipodal map of $S^3$ sends $\Omega$ to $-\Omega$ (\lean{oq\_antipodal}).

\section{Linear rigidity}

Here is the theorem that the author submitted for publication.

\begin{theorem}[linear rigidity]\proved\label{thm:rigidity}
Every linear perturbation of the Einstein--Yang--Mills Tolman universe that keeps spherical symmetry, the purely magnetic ansatz and conformal flatness is generated by a conformal Killing vector of the Einstein static universe. In particular it is pure gauge.
\end{theorem}
\noindent In the repository: the three generating modes \lean{rr\_mode0}, \lean{rr\_mode1}, \lean{rr\_mode2}, the integrability condition \lean{rr\_eta\_integrability}, the reduced equation \lean{rr\_chi\_ode} and the exclusion of the remaining mode \lean{rr\_cos\_mode\_excluded}.

The proof is local and uses neither mode expansions nor boundary conditions. The perturbations of $\Omega$ and $W$ must satisfy the linearised constraints; integrating them in $\eta$ and $\chi$ leaves a three-parameter family, and the three parameters are exactly the three conformal Killing vectors $\partial_\eta$, $\xi_1$, $\xi_2$ of the static universe that preserve spherical symmetry. They close into the algebra $\mathfrak{sl}(2,\mathbb R)$ (\lean{oq\_sl2\_bracket}), and the neck is the unique configuration invariant under all three (\lean{oq\_neck\_invariant}, \lean{oq\_neck\_unique}).

This does not contradict the known instabilities of the Einstein static universe or of Yang--Mills fields on a fixed background, because those perturbations leave the conformally flat class. The rigidity theorem says that within this class there is nothing to perturb.

\section{The self-dual connection}

Chapter~4 split Lorentz transformations into rotation and boost. General relativity has a matching split: the self-dual connection $A=\Gamma+iK$ combines the spin connection $\Gamma$ of space with its extrinsic curvature $K$. On the cap its coefficient is $c=1+i\dot a$, and the curvature coefficient $c^2-2c=-|c|^2$ is real. The Friedmann equation becomes $|c|^2=(8\pi G/3)\rho a^2$, and the phase of $c$ is the conformal time (\lean{curvature\_real}, \lean{friedmann\_selfdual}, \lean{phase\_lemma}). On both exact solutions the Yang--Mills field equals the real part of $A$ (\lean{cap\_real\_part}, \lean{real\_part}).

\section{A tempting quantisation, and why it was withdrawn}

The cap's action in the conformal frame is $S=(3\pi/4G)\int(\Omega^2-\dot\Omega^2)\,d\eta$, a harmonic oscillator of period $2\pi$ (\lean{periodic}, \lean{energy\_balance}). Quantising it naively gives $\ell^2/\ell_P^2=4(n+\tfrac12)/(3\pi)$ (\lean{quantized}), and with it a candidate value of the coupling. This was later withdrawn: in the full conformally flat magnetic sector the constraints leave no local degree of freedom, the quantum theory of that sector has a single state, and the coupling is not quantised (Chapter~23). \negative

\begin{standsbox}
The conformal transformation formula is quoted. The rigidity theorem is claimed as new (claim 1 of the guide for the specialist); a manuscript containing it was submitted to a journal and returned without review. The self-dual description is checked on the exact solutions only; in general evolution the statement ``Yang--Mills equals the real part of $A$'' is a postulate, not a theorem.
\end{standsbox}

\section*{Exercises}
\begin{exercise}Show that $\ell\cos\eta$ and $\ell/\sin u$ satisfy $\hat\Box\Omega=\Omega$. \emph{Answer:} \lean{cap\_box}, \lean{neck\_box}.\end{exercise}
\begin{exercise}Compute the three brackets of the conformal Killing vectors. \emph{Answer:} \lean{oq\_sl2\_bracket}.\end{exercise}
\begin{exercise}Show that $(1-x^2)U''-5xU'+n(n+4)U=0$ implies the Yang--Mills test-field frequencies $\omega^2=(n+2)^2-3$. \emph{Answer:} \lean{oq\_spectrum}.\end{exercise}
